% Volume IV: Law as an Epistemic Machine
% Part VII. Evidence Before Action
% Chapter 36: Burdens and Standards of Proof
% Spine: proof-spines/ch036-spine.md

\chapter{Burdens and Standards of Proof}
\label{ch:ch036}

Burdens and standards of proof compare probable cause, reasonable suspicion, preponderance of evidence, clear and convincing evidence, and proof beyond a reasonable doubt as different permissions to act under uncertainty. They are not merely different numerical thresholds, but an approximate threshold model clarifies their function:
\[
A\text{ permitted if }P(H\mid E)>\theta_A.
\]
\ToyModel\ The point is not that every doctrine secretly encodes a precise scalar \(\theta_A\).
It is that each standard specifies how much evidentiary support must exist before some action \(A\) becomes licit. As the magnitude, irreversibility, and asymmetry of the contemplated intervention rise, the threshold should rise with them. Standards of proof are therefore calibrations of authorized action under uncertainty rather than meters of metaphysical certainty. And when a standard demands evidentiary movement no institution can realistically supply for the kind of harm at issue, the same rigor flips into an unmeetable burden that protects the uncertainty-producing status quo.
